% === Introduction and motivation ===

The document describes the algorithms, data structures, and implementation
of \FractalShark{}, a high-performance, GPU-accelerated Mandelbrot set
renderer designed for extreme deep-zoom exploration.  \FractalShark{}
combines perturbation theory with linear approximation, NTT-based
high-precision GPU arithmetic, and optimized reference-orbit computation to
enable interactive navigation at zoom depths exceeding $10^{10000}$.  Its
Feature Finder automatically locates high-period periodic points, enabling
targeted exploration of deeply embedded structures without manual search.

The project has evolved through extensive experimentation with numerical
precision, memory management, and GPU kernel design.  As a result, this
document serves a dual purpose: it is both \emph{technical documentation}
of the algorithms and architecture underlying \FractalShark{}, and a
\emph{research notebook} recording the design rationale, tradeoffs, and
lessons learned during development.  Where relevant, alternative
approaches that were considered or attempted are noted, along with
explanations of why the current design was chosen.

The presentation proceeds from high-level goals and mathematical background
through the individual subsystems---rendering kernels, perturbation
pipelines, reference-orbit computation, GPU arithmetic, and the Feature
Finder---and concludes with infrastructure details and development history.

\section{\FractalShark{} Overview}
\label{sec:goal}

\paragraph{Reader background and route through the notes.}
The mathematical explanations assume elementary calculus, complex arithmetic, and
basic programming. Numerical methods, transform multiplication, and GPU execution
concepts are introduced before their implementation details. The rendering chapters
first establish a direct calculation, then show which work perturbation and
approximation can share or skip. The reference-orbit chapters explain how that shared
input is computed. Readers interested in using the application can begin with
\cref{sec:user-documentation}; the engineering appendix and dated development history
provide implementation and historical context. The final appendix treats GPU PNG
encoding as a supporting technique rather than a prerequisite for fractal mathematics.

\begin{figure}[!htbp]
  \centering
  \includegraphics[width=0.9\linewidth]{mandelbrot.png}
  \caption[The Mandelbrot set]{The Mandelbrot set.}
  \label{fig:mandelbrot}
\end{figure}

\FractalShark{} is an interactive Mandelbrot set renderer (\cref{fig:mandelbrot}) focused on
\textbf{extreme deep-zoom exploration}, numerical correctness, and
algorithmic experimentation.  It supports zoom depths ranging from
conventional views to magnifications exceeding $10^{10000}$, while
exposing internal rendering choices and precision tradeoffs to the user.

\paragraph{Interactive navigation and view control.}
\FractalShark{} provides direct mouse-driven navigation, including centering
the view at an arbitrary point, zooming in and out at the cursor location,
stepping backward through navigation history, and invoking automatic zoom
modes.  Window management options allow toggling between windowed and
full-screen modes, including square-aspect rendering for precise analysis.

\paragraph{Built-in views for demonstration and validation.}
The application includes a set of built-in views covering a wide
range of numerical regimes: GPU precision limits, high-period locations,
known hard points, historical bug cases, regression tests, and ultra-deep
zoom demonstrations.  These views serve both as showcases of capability and
as repeatable test cases for performance and correctness.

\paragraph{Explicit algorithm selection and transparency.}
A central design goal of \FractalShark{} is to make rendering algorithms
explicit and selectable.  Users may rely on an automatic algorithm selector
or manually choose among many rendering paths (see \cref{subsec:ref-orbit-mt-tradeoff} for automatic selection), including:
\begin{itemize}
  \item GPU-based low-zoom renderers with selectable iteration precision and
        multiple numeric formats (single-, dual-, and quad-limb 32- and 64-bit
        variants, as well as deeper representations discussed subsequently);
        see \cref{sec:mandel-base-kernels}.
  \item Scaled perturbation algorithms for extending low-precision arithmetic
        to deeper zooms (\cref{sec:perturbation-concept}).
  \item The retained bilinear approximation (BLA~v1) perturbation paths
        (\cref{sec:bilinear-approx}).
  \item The current linear approximation (LA~v2) pipelines, including full rendering,
        linear-approximation-only, and perturbation-only variants
        (\cref{sec:la-v2-perturb}).
  \item Reference-compression–aware algorithms, including Imagina-compatible
        ``max'' compression formats (\cref{subsec:ref-orbit-compression-reuse}).
  \item CPU-only algorithms for very high precision, verification, and
        fallback scenarios.
\end{itemize}

Many of these modes are exposed explicitly for testing and comparison, and
not all are intended as default or production-quality paths.

LA~v2 is the primary deep-zoom approximation path. Its nonlinear
input preparation and adaptive stage boundaries address limitations
of BLA~v1 (\cref{sec:la-v2-motivation}). Automatic selection can still
use other algorithms at shallower zooms, including BLA~v1 in the CPU
fallback tiers.

\paragraph{Linear approximation configuration.}
Linear approximation behavior can be configured for multithreaded or
single-threaded execution and tuned via presets that prioritize accuracy,
performance, or memory usage.  These controls reflect ongoing development
and experimentation with LA parameter tradeoffs.

\paragraph{Image quality and coloring.}
\FractalShark{} supports GPU antialiasing at multiple sample levels and a
flexible palette system (\cref{subsec:antialiasing,subsec:coloring}).  Palettes may be selected from predefined themes,
generated randomly, and rendered at configurable color depths.  Some palette
features are intentionally limited or disabled where they are known to be
incomplete.

\paragraph{Iteration limits and precision control.}
Users can dynamically adjust iteration limits, switch between 32-bit and
64-bit iteration counters, and modify shallow-zoom iteration precision.
These controls are primarily relevant for low-zoom or testing scenarios and
are not universally applicable to all rendering algorithms.

\paragraph{Perturbation and reference-orbit management.}
\FractalShark{} places strong emphasis on reference-orbit reuse and
perturbation-based rendering (\cref{sec:ref-orbit-calc}).  It supports multiple perturbation strategies
(single-threaded, multithreaded, periodicity-assisted (\cref{sec:hp-periodicity}), and GPU
variants), along with tools to clear, inspect, save, reload, and reuse
reference orbits.  Automatic and manual reference-orbit persistence is
supported via file-backed storage (\cref{sec:disk_backed_growable_vectors}).

\paragraph{Memory management and scalability.}
To enable extreme zoom depths without exhausting system memory, \FractalShark{}
employs file-backed storage, optional reference compression, configurable
memory limits, and automatic cleanup policies.  These mechanisms are closely
tied to recent allocator and reference-orbit refactors discussed in the
project development history.

\paragraph{Diagnostics, benchmarking, and testing.}
The application includes facilities for displaying detailed rendering
parameters, running repeatable benchmarks, executing regression tests, and
comparing reference orbits (including Imagina-compatible formats).  These
features reflect \FractalShark{}'s dual role as both a visualization tool and a
development platform.

\paragraph{Data export and interoperability.}
\FractalShark{} can save rendered images, high-resolution bitmaps, iteration
counts, and reference orbits in multiple text and compressed formats.  It
supports loading external reference orbits and matching Imagina-compatible
data for cross-tool validation.

\medskip
\noindent
Overall, \FractalShark{} functions as both a \textbf{deep-zoom Mandelbrot
explorer} and a \textbf{research and experimentation platform} for
high-precision fractal rendering, prioritizing transparency, correctness, and
performance.  It is not intended as a polished end-user application, but rather
as an experimental platform for exploring and validating advanced fractal
rendering algorithms.
